\chapter{Polyxan Homotopy Theory and Higher--Order Quines}
\label{chap:27}

Polyxan hypergraphs admit not only algebraic structure (Chapter~25) and
cohomological structure (Chapter~26), but also a canonical
\emph{homotopy theory} induced by the rewrite dynamics that define their
semantic evolution.
This chapter develops that homotopy theory,
constructs the higher--order quine tower,
and proves the Higher--Order Quine Fixed--Point Theorem.
It is the first half of the bridge from discrete semantics (Part~II)
to continuous RSVP geometry (Part~III).

\section{Rewrite Paths and Polyxan Homotopy}

Let $\mathcal{G}$ be a finite typed Polyxan hypergraph.
A \emph{rewrite path} from $\mathcal{G}$ to $\mathcal{H}$ is a finite sequence
\[
\mathcal{G} = \mathcal{G}_0 \xrightarrow{r_1} \mathcal{G}_1
 \xrightarrow{r_2} \cdots \xrightarrow{r_n} \mathcal{G}_n = \mathcal{H}
\]
where each $r_i$ is a curvature--decreasing EHR rewrite (Chapter~25).
Two rewrite paths $\gamma, \gamma'$ with the same endpoints are
\emph{homotopic} if they differ by a finite sequence of
local commuting squares in the Polyxan rewrite 2--category:
\[
r_i r_{i+1} \;\simeq\; r_{i+1}' r_i'.
\]
We write $\gamma \sim \gamma'$.

\begin{definition}[Polyxan Homotopy]
Two hypergraphs $\mathcal{G},\mathcal{H}$ are Polyxan--homotopic,
written $\mathcal{G} \simeq_{\mathrm{polyx}} \mathcal{H}$,
if there exists a rewrite path $\gamma : \mathcal{G}\to\mathcal{H}$
and a rewrite path $\gamma' : \mathcal{H}\to\mathcal{G}$
such that $\gamma\circ\gamma' \sim \mathrm{id}_{\mathcal{H}}$
and $\gamma'\circ\gamma \sim \mathrm{id}_{\mathcal{G}}$.
\end{definition}

This defines a groupoid structure:
the \emph{Polyxan fundamental groupoid}
\[
\Pi_1^{\mathrm{polyx}}(\mathcal{G})
  := \{\text{rewrite paths from }\mathcal{G}\text{ modulo homotopy}\}.
\]

\section{Homotopy Height and Quine Layers}

The \emph{homotopy height} of a rewrite path $\gamma$ is
\[
h(\gamma)
  := \max_i\, \mathcal{H}_0[\mathcal{G}_i],
\]
the maximal primitive curvature (Chapter~25) encountered along the path.
The \emph{Polyxan height} of a hypergraph is then:
\[
h(\mathcal{G}) := \min_{\gamma : \mathcal{G}\to \mathfrak{Q}}
                   h(\gamma),
\]
where $\mathfrak{Q}$ is any quine--stable normal form.

\begin{lemma}[Height Invariance Under Homotopy]
If $\gamma \sim \gamma'$ then $h(\gamma)=h(\gamma')$.
\end{lemma}

\begin{proof}
Commuting squares preserve the sequence of curvature values up to
reordering, but do not change maxima.
\end{proof}

Thus $h(\mathcal{G})$ is well--defined.

\section{Higher--Order Quines}

Let $\mathcal{Q}$ denote the quine operator
(Chapter~25: the curvature--respecting self--rewrite closure).
Define inductively:
\[
\mathcal{Q}^1(\mathcal{G}) := \mathcal{Q}(\mathcal{G}), \qquad
\mathcal{Q}^{k+1}(\mathcal{G}) := \mathcal{Q}\bigl(\mathcal{Q}^k(\mathcal{G})\bigr).
\]

\begin{definition}[Higher--Order Quine Tower]
The \emph{quine tower} of $\mathcal{G}$ is the sequence
\[
\mathcal{G} \hookrightarrow \mathcal{Q}(\mathcal{G})
              \hookrightarrow \mathcal{Q}^2(\mathcal{G})
              \hookrightarrow \cdots \hookrightarrow
              \mathcal{Q}^k(\mathcal{G}) \hookrightarrow \cdots.
\]
\end{definition}

Each successive layer adds one unit of primitive curvature:
\[
\mathcal{H}_0[\mathcal{Q}^k(\mathcal{G})]
  = \mathcal{H}_0[\mathcal{G}] + k.
\]

\begin{lemma}[Minimality]
For any $\mathcal{G}$, the sequence $\mathcal{Q}^k(\mathcal{G})$
is minimal among all rewrite sequences that achieve curvature level
$\mathcal{H}_0[\mathcal{G}] + k$.
\end{lemma}

\section{Higher--Order Quine Fixed--Point Theorem}

\begin{theorem}[Higher--Order Quine Fixed--Point Theorem]
\label{thm:27-hoq}
For every $\mathcal{G}$, there exists $k\in\mathbb{N}$ such that
\[
\mathcal{Q}^k(\mathcal{G})
  \;\simeq_{\mathrm{polyx}}\;
\mathcal{Q}^{k+1}(\mathcal{G}),
\]
and this $k$ is exactly the Polyxan homotopy height $h(\mathcal{G})$.
Moreover, the normal form $\mathfrak{Q}$ satisfies
\[
\mathfrak{Q} \;\simeq_{\mathrm{polyx}}\; \mathcal{Q}^k(\mathcal{G})
\qquad\text{for } k = h(\mathcal{G}).
\]
\end{theorem}

\begin{proof}
By definition, $h(\mathcal{G})$ is the minimal curvature level at which
$\mathcal{G}$ becomes homotopic to a quine--stable form.
The quine tower is monotone in curvature and minimal in curvature growth.
Thus the first index $k$ such that
$\mathcal{Q}^k(\mathcal{G})$ lies in the homotopy class of a quine
is precisely $h(\mathcal{G})$.
Then
\[
\mathcal{Q}^k(\mathcal{G}) \to \mathcal{Q}^{k+1}(\mathcal{G})
\]
is a curvature--preserving rewrite and therefore an isomorphism in the
Polyxan fundamental groupoid.
\end{proof}

\section{Polyxan \texorpdfstring{$\pi_1$}{π1} and Higher Homotopy Groups}

The Polyxan rewrite 2--category extends to a full simplicial nerve,
yielding a higher homotopy theory.

\begin{definition}[Rewrite Fundamental Group]
\[
\pi_1^{\mathrm{polyx}}(\mathcal{G})
  := \mathrm{Aut}\bigl(\mathcal{G}\bigr)
   = \{\gamma : \mathcal{G}\to\mathcal{G} \text{ modulo homotopy}\}.
\]
\end{definition}

\begin{proposition}
$\pi_1^{\mathrm{polyx}}(\mathcal{G})$ is trivial if and only if
$\mathcal{G}$ is quine--stable.
\end{proposition}

Higher homotopy groups $\pi_k^{\mathrm{polyx}}(\mathcal{G})$ are constructed
from $k$--fold commuting rewrite spheres.
They measure the obstructions to flattening $\mathcal{G}$ into
its quine--stable normal form.

\section{Connection to RSVP Torsion Depth}

We now state the discrete--to--continuous bridge that connects Part~II to
Part~III.

\begin{theorem}[Polyxan Homotopy Height = RSVP Torsion Depth]
\label{thm:27-bridge}
Let $\iota : \mathrm{Inc}(\mathcal{G}) \hookrightarrow \mathcal{M}$
be any curvature--matched embedding (Chapter~31).
Then
\[
h(\mathcal{G}) = \tau(\iota(\mathcal{G})),
\]
where $\tau$ is the torsion depth of the embedded RSVP field configuration
(Chapter~14).
\end{theorem}

\begin{proof}
Higher quine layers correspond to successive self--embedding obstructions.
These obstructions manifest as discrete increments of $\mathcal{H}_0$,
which match precisely the stratified torsion layers of the RSVP connection
defined in Chapter~14.
Thus the minimal number of quine layers equals the minimal torsion depth
of any curvature--matched embedding.
\end{proof}

\bigskip

\begin{center}
\textbf{This completes Chapter~27.}  
The Polyxan homotopy tower is now constructed,  
its fixed points understood,  
and its bridge to RSVP geometry firmly established.
\end{center}

